\section{Lecture 9: Feature Maps and Nonlinear Models}

\subsection{Motivation}

In previous lectures, we studied linear models such as linear and logistic regression.

\medskip

\noindent
While these models are simple and interpretable, they are limited:

\begin{itemize}
    \item They can only represent linear relationships
    \item Many real-world patterns are inherently nonlinear
\end{itemize}

\medskip

\noindent
\textbf{Key question.}  
How can we model nonlinear relationships while retaining the simplicity of linear models?

\medskip

\noindent
\textbf{Guiding idea.}  
Transform the input representation so that nonlinear patterns become linear.

\subsection{Basis Expansions}

We introduce a feature map:
\[
\phi : \mathcal{X} \to \mathbb{R}^D,
\]
which transforms inputs into a higher-dimensional feature space.

\medskip

\noindent
\textbf{Examples.}

\begin{itemize}
    \item Polynomial features:
    \[
    \phi(x) = (1, x, x^2, \dots, x^d)
    \]

    \item Interaction terms:
    \[
    \phi(x_1, x_2) = (x_1, x_2, x_1 x_2)
    \]
\end{itemize}

\medskip

\noindent
\textbf{Model.}
\[
f(x) = w^\top \phi(x).
\]

\medskip

\noindent
\textbf{Interpretation.}
\begin{itemize}
    \item Model is linear in $\phi(x)$
    \item May be nonlinear in original input $x$
\end{itemize}

\medskip

\noindent
\textbf{Insight.}  
Nonlinearity can be introduced through feature transformations.

\subsection{Linear Models in Feature Space}

We now apply linear regression or logistic regression to $\phi(x)$:

\[
f(x) = w^\top \phi(x).
\]

\medskip

\noindent
\textbf{Key observation.}
\begin{itemize}
    \item The learning algorithm remains unchanged
    \item Only the representation changes
\end{itemize}

\medskip

\noindent
\textbf{Decision boundaries.}
\begin{itemize}
    \item Linear in feature space
    \item Nonlinear in input space
\end{itemize}

\medskip

\noindent
\textbf{Example.}
\begin{itemize}
    \item A quadratic feature map produces curved decision boundaries
\end{itemize}

\medskip

\noindent
\textbf{Insight.}  
Feature maps allow us to trade simplicity of models for richness of representation.

\subsection{The Kernel Idea (Intuition)}

Working with high-dimensional feature maps can be computationally expensive.

\medskip

\noindent
\textbf{Key observation.}  
Many learning algorithms depend only on inner products:
\[
\langle \phi(x), \phi(x') \rangle.
\]

\medskip

\noindent
\textbf{Kernel function.}
\[
k(x, x') = \langle \phi(x), \phi(x') \rangle.
\]

\medskip

\noindent
\textbf{Examples.}
\begin{itemize}
    \item Polynomial kernel:
    \[
    k(x, x') = (x^\top x' + c)^d
    \]

    \item Gaussian kernel:
    \[
    k(x, x') = \exp(-\gamma \|x - x'\|^2)
    \]
\end{itemize}

\medskip

\noindent
\textbf{Kernel trick.}
\begin{itemize}
    \item Compute inner products without explicitly computing $\phi(x)$
    \item Enables work in very high-dimensional spaces
\end{itemize}

\medskip

\noindent
\textbf{Insight.}  
We can operate in rich feature spaces implicitly through kernels.

\subsection{Expressivity through Features}

The expressive power of a model depends heavily on the choice of features.

\medskip

\noindent
\textbf{Key idea.}
\begin{itemize}
    \item A linear model with rich features can approximate complex functions
\end{itemize}

\medskip

\noindent
\textbf{Tradeoffs.}
\begin{itemize}
    \item More expressive features:
    \begin{itemize}
        \item Lower bias
        \item Higher risk of overfitting
    \end{itemize}
\end{itemize}

\medskip

\noindent
\textbf{Connection to inductive bias.}
\begin{itemize}
    \item Feature map defines the hypothesis space
\end{itemize}

\medskip

\noindent
\textbf{Modern perspective.}
\begin{itemize}
    \item Deep learning learns feature maps automatically
\end{itemize}

\medskip

\noindent
\textbf{Insight.}  
Representation is a central driver of model performance.

\subsection{A Unifying Perspective}

Feature maps unify linear and nonlinear modeling:

\begin{itemize}
    \item \textbf{Linear model:} simple structure
    \item \textbf{Feature map:} expressive representation
    \item \textbf{Kernel:} implicit computation in feature space
\end{itemize}

\medskip

\noindent
\textbf{Core idea.}  
Complex learning problems can be solved by combining simple models with powerful representations.

\medskip

\noindent
\textbf{Key insight.}  
The boundary between linear and nonlinear models is often a matter of representation.

\subsection{Outlook}

In this lecture, we introduced:

\begin{itemize}
    \item basis expansions and feature maps,
    \item linear models in feature space,
    \item the kernel idea,
    \item expressivity through representations.
\end{itemize}

\medskip

\noindent
These ideas lead naturally to kernel-based methods.

\medskip

\noindent
In the next lecture, we study support vector machines, where kernels play a central role.

\medskip

\noindent
\textbf{Guiding principle.}  
Representation transforms the problem; learning solves it.